\label{der:8-4}
Jordan形的存在与唯一已在推导\ref{der:jordan-full}由链基完整证明。现逐块分析。令$J_s(\lambda)=\lambda I+N_s$，$N_s$把第$j$个基向量送到前一向量，首向量送零，因此$\rank N_s^k=\max(s-k,0)$、$\dim\ker N_s^k=\min(k,s)$。若块属于另一根$\mu\ne\lambda$，$J_s(\mu)-\lambda I$为对角元全$\mu-\lambda\ne0$的三角矩阵，幂仍可逆，不贡献核。对所有属于$\lambda$的块求和便得$d_k$。相邻差$\min(k,s)-\min(k-1,s)$在$s\ge k$时为1，否则为0，所以至少$k$的块数为$d_k-d_{k-1}$；再减去至少$k+1$的块数，得到恰为$k$的块数$2d_k-d_{k-1}-d_{k+1}$。

把多项式在$\lambda$附近作有限展开$p(t)=\sum_{j\ge0}b_j(t-\lambda)^j$，系数$b_j=\frac{p^{(j)}(\lambda)}{j!}$可由多项式逐项求导验证。代入块后
\[p(J_s(\lambda))=b_0I+b_1N_s+\cdots+b_{s-1}N_s^{s-1}.\]
这些矩阵的非零位置分别在主对角、第1至$s-1$超对角，互不重叠，因此和为零当且仅当$b_0=\cdots=b_{s-1}=0$，也就是$p$含$(t-\lambda)^s$。对全部块同时零化，必须对每个根取最大块长度，乘起来便是最低次首一零化多项式。

所有块长为1恰当于Jordan形已对角，也恰当于最小多项式每个根指数为1。若数域不分裂，即使无重因子也未必能在该数域对角化，所以“分裂”条件仍须保留。
